%%
%% Speech Analytics and Natural Language Processing Applied to
%% Clinical Communication: A Scoping Review
%%
%% Rafael Baena Neto · Vicente Idalberto Becerra Sablón
%% Graduate Program in Health Data Science — Universidade São Francisco (USF)
%% Target journal: International Journal of Medical Informatics (IJMI)
%% APC covered: CAPES/Elsevier Read & Publish 2026–2028
%%
%% Compiled with: cas-sc.cls v2.4 (single-column, for submission/review)
%% Dependencies (copy from /latex/): cas-sc.cls, cas-common.sty,
%%   cas-model2-names.bst
%%

\documentclass[a4paper,fleqn]{cas-sc}

\usepackage[authoryear,longnamesfirst]{natbib}
\usepackage[utf8]{inputenc}
\usepackage[T1]{fontenc}
\usepackage{lmodern}
\usepackage{amsmath}
\usepackage{graphicx}
\usepackage{booktabs}
\usepackage{longtable}
\usepackage{array}
\usepackage{hyperref}
\usepackage{xcolor}

\hypersetup{colorlinks=true, linkcolor=blue, citecolor=blue, urlcolor=blue}

%% ================================================================
\begin{document}
\let\WriteBookmarks\relax
\def\floatpagepagefraction{1}
\def\textpagefraction{.001}

%% ----------------------------------------------------------------
%% Running head
%% ----------------------------------------------------------------
\shorttitle{Speech Analytics in Clinical Communication}
\shortauthors{R. Baena Neto et~al.}

%% ----------------------------------------------------------------
%% Title
%% ----------------------------------------------------------------
\title[mode=title]{Speech Analytics and Natural Language Processing
Applied to Clinical Communication: A Scoping Review}

%% ----------------------------------------------------------------
%% Authors
%% Affiliation [1]: Graduate Program in Health Data Science, USF
%% Affiliation [2]: Department of Computer Engineering, USF (IC students)
%% Rafael  ORCID: 0009-0005-2122-4665
%% Vicente ORCID: 0000-0003-3127-1906
%% IC students ORCID: register at orcid.org before submission
%% ----------------------------------------------------------------
\author[1]{Rafael Baena Neto}[%
  type=editor,
  orcid=0009-0005-2122-4665,
  role=Researcher
]
\cormark[1]
\fnmark[1]
\ead{rafael.baena@mail.usf.edu.br}

\credit{Conceptualisation, Methodology, Data curation, Formal analysis,
        Writing -- original draft}

\affiliation[1]{%
  organization={Graduate Program in Health Data Science,
                Universidade São Francisco (USF)},
  addressline={Rua Alexandre Rodrigues Barbosa, 45},
  city={Bragança Paulista},
  postcode={12916-900},
  state={SP},
  country={Brazil}
}

\author[1]{Vicente Idalberto Becerra Sablón}[%
  orcid=0000-0003-3127-1906
]
\fnmark[2]
\ead{vicente.sablon@usf.edu.br}

\credit{Supervision, Conceptualisation, Writing -- review and editing}

\author[2]{Fabrício Evangelista Dos Santos}[%
  orcid=0000-0000-0000-0000
]
\fnmark[3]
\ead{fabricio.evangelista@mail.usf.edu.br}

\credit{Data curation, Formal analysis, Writing -- review and editing}

\author[2]{Lucas Kenzo Terazzin Ida}[%
  orcid=0000-0000-0000-0000
]
\fnmark[4]
\ead{lucas.ida@mail.usf.edu.br}

\credit{Data curation, Formal analysis, Writing -- review and editing}

\author[2]{Murilo Nogueira Minutti}[%
  orcid=0000-0000-0000-0000
]
\fnmark[5]
\ead{murilo.minutti@mail.usf.edu.br}

\credit{Data curation, Formal analysis, Writing -- review and editing}

\affiliation[2]{%
  organization={Department of Computer Engineering,
                Universidade São Francisco (USF)},
  addressline={Rua Alexandre Rodrigues Barbosa, 45},
  city={Bragança Paulista},
  postcode={12916-900},
  state={SP},
  country={Brazil}
}

%% ----------------------------------------------------------------
%% Footnotes / correspondence
%% ----------------------------------------------------------------
\cortext[1]{Corresponding author}
\fntext[1]{}
\fntext[2]{}
\fntext[3]{}
\fntext[4]{}
\fntext[5]{}

%% ----------------------------------------------------------------
%% Abstract (250 words max — IJMI)
%% Current count: ~235 words
%% ----------------------------------------------------------------
\begin{abstract}
\textbf{Background:}
Professional–patient communication is a dense yet predominantly unstructured
and analytically underutilised source of clinical information. Advances in
automatic speech recognition (ASR), natural language processing (NLP), and
vocal biomarker analysis have expanded the capacity to extract structured data
from clinical interactions; nonetheless, few studies operationalise these
interactions as measurable speech-derived variables, particularly
communicational and behavioural markers relevant to patient assessment.

\textbf{Objective:}
To identify, synthesise, and critically evaluate evidence on the application
of speech analytics and NLP to professional–patient clinical communication,
mapping methodological approaches, identified markers, technical challenges,
and research gaps.

\textbf{Methods:}
Scoping review following PRISMA-ScR \citep{Tricco2018} and JBI guidelines
\citep{Peters2020}. The corpus comprised 27 articles selected from 36 records
retrieved from three documentary bases, organised across four thematic axes:
(A)~clinical NLP and biomedical entity extraction; (B)~automated clinical
documentation and ASR; (C)~physician–patient communication; (D)~vocal
biomarkers and multimodal analysis.

\textbf{Results:}
Included studies document advances in clinical NLP (transformer-based models,
BERT leading named-entity recognition) and automated transcription (word error
rate $\approx$9\% with Whisper/PyAnnote). Patient-centred communication
demonstrates measurable impact on objective health outcomes. Communicational
markers — self-repair, engagement, emotional support, and barrier indicators —
are identified as adherence predictors. Multimodal acoustic–semantic analysis
remains nascent; most approaches focus on documentary automation rather than
treating clinical interaction as an analytic object.

\textbf{Conclusions:}
A consistent gap is confirmed: few studies operationalise clinical interactions
as measurable speech-derived variables, particularly communicational and
behavioural markers extracted from anamnesis. This review supports the design
of prospective studies integrating speech analytics into structured clinical
communication analysis for adherence monitoring and care quality assessment.

\textbf{Registration:} PROSPERO [CRD number to be assigned]
\end{abstract}

%% ----------------------------------------------------------------
%% Research Highlights — MANDATORY for IJMI (3 bullet points)
%% Each item: ≤ 85 characters
%% ----------------------------------------------------------------
\begin{highlights}
\item BERT-based transformer models lead clinical NER; WER $\approx$9\%
      achievable with Whisper/PyAnnote pipelines
\item Communicational markers (self-repair, engagement, barrier indicators)
      predict therapeutic adherence
\item Consistent gap: clinical interactions rarely operationalised as
      measurable speech-derived variables
\end{highlights}

%% ----------------------------------------------------------------
%% Keywords (4–6 per IJMI guidelines)
%% ----------------------------------------------------------------
\begin{keywords}
Speech analytics \sep
Natural language processing \sep
Clinical communication \sep
Automatic speech recognition \sep
Communicational markers \sep
Digital scribe
\end{keywords}

\maketitle

%% ================================================================
\section{Introduction}
\label{sec:intro}

Clinical communication constitutes one of the richest yet most underutilised
sources of health data. In public health systems, anamneses and therapeutic
orientations predominantly occur as unstructured verbal exchanges rarely
converted into analysable data \citep{Wang2018,Koleck2019}.

The quality of the professional–patient interaction has been shown to
influence adherence and objective clinical outcomes \citep{Kelley2014}.
Communication barriers — including health literacy gaps, misaligned
expectations, and miscommunication patterns — are associated with impaired
therapeutic adherence and poorer prognosis
\citep{McCabe2018,Alkhamees2023}.

Speech analytics — the computational analysis of spoken language — offers a
pathway to extract structured, measurable indicators from verbal clinical
interactions. Recent advances in ASR, particularly transformer-based models
such as Whisper \citep{Radford2023}, alongside instruction-tuned large
language models (LLMs), have lowered barriers to deploying end-to-end NLP
pipelines in healthcare \citep{Klusty2025}. Nonetheless, most existing
approaches target diagnostic support or documentation automation rather than
treating the clinical interaction itself as an object of structured analysis
\citep{vanBuchem2021}. The development of a speech analytics protocol for
monitoring orthopedic rehabilitation in a public health system
\citep{BaenaNeto2026} exemplifies the translational potential of these
technologies and underscores the need for a consolidated evidence map to
support prospective study design.

This scoping review aims to map the evidence on speech analytics and NLP
applied to clinical communication, characterising the methodological
landscape and identifying gaps for future prospective studies.

%% ================================================================
\section{Methods}
\label{sec:methods}

\subsection{Study design}

This scoping review follows the methodological framework proposed by
\citet{Arksey2005}, refined by \citet{Levac2010} and the Joanna Briggs
Institute (JBI) \citep{Peters2020}. Reporting adheres to the Preferred
Reporting Items for Systematic Reviews and Meta-Analyses extension for
Scoping Reviews (PRISMA-ScR) \citep{Tricco2018}.

\subsection{Research questions}

\begin{itemize}
  \item \textbf{Q1:} What speech analytics approaches have been applied to
        clinical communication in the last five years?
  \item \textbf{Q2:} What NLP and ASR methods are used for automated
        transcription and analysis of clinical consultations?
  \item \textbf{Q3:} What communicational and behavioural markers have been
        identified or proposed in professional–patient interaction?
  \item \textbf{Q4:} What are the technical, ethical, and methodological
        challenges for automated analysis of clinical communication?
\end{itemize}

\subsection{Eligibility criteria}

\subsubsection{Inclusion criteria}

\begin{itemize}
  \item \textbf{IC1:} Study applies speech analytics, NLP, or ASR to a
        clinical context
  \item \textbf{IC2:} Full text available in English, Portuguese, or Spanish
  \item \textbf{IC3:} Analyses professional–patient communication or
        information extraction from clinical consultations
  \item \textbf{IC4:} Describes methods, results, or markers applicable to
        clinical interaction
  \item \textbf{IC5:} Published between January 2014 and 2026
\end{itemize}

\subsubsection{Exclusion criteria}

\begin{itemize}
  \item \textbf{EC1:} Publications prior to 2014
  \item \textbf{EC2:} Extended abstracts, tutorials, or white papers without
        empirical data
  \item \textbf{EC3:} Exclusively focused on voice-based diagnosis of physical
        pathologies without analysis of the communicational interaction
  \item \textbf{EC4:} Studies on non-clinical populations without explicitly
        stated methodological transferability
\end{itemize}

\subsection{Information sources and search strategy}

Searches were conducted in [specify the three documentary bases]. The
combined Boolean search string is presented below.

\begin{verbatim}
("speech analytics" OR "automatic speech recognition" OR "ASR"
OR "natural language processing" OR "NLP" OR "large language model"
OR "speech-to-text" OR "voice recognition" OR "vocal biomarkers"
OR "speaker diarization")
AND
("clinical communication" OR "doctor-patient" OR "physician-patient"
OR "anamnesis" OR "medical consultation" OR "health interaction"
OR "patient-clinician communication" OR "therapeutic adherence")
AND
("clinical notes" OR "SOAP notes" OR "medical transcription"
OR "digital scribe" OR "electronic health record"
OR "communicational markers" OR "behavioural markers")
\end{verbatim}

\subsection{Study selection}

Titles and abstracts were screened independently by two reviewers. A total
of 36 records were retrieved from three documentary bases. After two-phase
screening (title/abstract then full text), 27 articles were included in the
final corpus. Disagreements were resolved by consensus or adjudication by a
third reviewer. The selection process is documented in a PRISMA-ScR flow
diagram (Fig.~\ref{fig:prisma}).

\subsection{Data charting}

Data were extracted using a standardised charting form including: author(s),
year, country, study design, clinical setting, population characteristics,
technology stack (ASR model, NLP framework, LLM), clinical variables
extracted, outcomes reported, ethical framework, and limitations.

\subsection{Synthesis}

A narrative synthesis was conducted. Findings are organised around the
four thematic axes defined in eligibility screening: (A)~clinical NLP
and biomedical entity extraction, (B)~automated documentation and ASR,
(C)~physician–patient communication, and (D)~vocal biomarkers and
multimodal analysis.

%% ================================================================
\section{Results}
\label{sec:results}

%% [To be completed after full data extraction]
%%
%% Suggested structure:
%% 3.1 Study selection
%%     PRISMA-ScR flow diagram: 36 screened → 27 included
%%     Table of excluded studies with reasons
%% 3.2 Characteristics of included studies
%%     Table: code | authors | year | type | axis | main metric | notes
%% 3.3 Axis A — Clinical NLP and biomedical entity extraction
%%     Key findings: transformer transition, BERT NER leadership,
%%     BERT-CRF F1 > 0.8 (Watabe 2025), zero-shot NER (Sezgin 2023)
%% 3.4 Axis B — Automated documentation and ASR
%%     Key findings: WER ~9% (Klusty 2025), SOAP generation (Krishna 2021,
%%     Schloss 2020), diarisation challenges (Tran 2023)
%% 3.5 Axis C — Physician–patient communication
%%     Key findings: self-repair as adherence predictor (McCabe 2018),
%%     DPCQ scale (Shao 2025), meta-analysis on outcomes (Kelley 2014),
%%     health literacy (Muscat 2020), intercultural barriers (Alkhamees 2023)
%% 3.6 Axis D — Vocal biomarkers and multimodal analysis
%%     Key findings: nascent state (Albert 2024), acoustic pipeline
%%     (Fagherazzi 2021), topic modelling (Imel 2022)
%% 3.7 Ethical and regulatory frameworks
%%     Key findings: on-premise requirement (Klusty 2025), LLM safety
%%     metrics (Abbasian 2024)

%% ================================================================
\section{Discussion}
\label{sec:discussion}

%% [To be completed]
%%
%% Key themes:
%% - Methodological heterogeneity across four axes
%% - Gap confirmed: clinical interaction treated as means, not object
%% - Communicational and behavioural markers from anamnesis: underexplored
%% - Implications for prospective protocol design (Fase 2 / Sablon protocol)
%% - Limitations: English-dominant corpus, three-base coverage,
%%   heterogeneous study quality, no date restriction pre-2014
%% - Ethical note: voice as biometry, LGPD, consent requirements,
%%   on-premise processing

%% ================================================================
\section{Conclusion}
\label{sec:conclusion}

%% [To be completed]
%% Core points:
%% - Consolidated advances in clinical NLP and ASR confirmed
%% - Consistent gap: few studies operationalise clinical interactions as
%%   measurable speech-derived variables
%% - Need for prospective studies integrating speech analytics into
%%   structured clinical communication analysis

%% ================================================================
\section*{Declarations}

\subsection*{Funding}
This study received no specific funding.
The article processing charge (APC) is covered under the Read and Publish
agreement between CAPES and Elsevier (2026--2028).

\subsection*{Conflict of interest}
The authors declare no conflict of interest.

\subsection*{Ethics approval}
Not applicable. This scoping review is based on publicly available
published literature and does not involve human participants.

\subsection*{Data availability}
The data charting form and full search strategies are available as
supplementary material upon acceptance.

\subsection*{AI use disclosure}
In accordance with Elsevier's Generative AI Policy: AI tools were used
solely for language support and literature searches. All intellectual
content, analysis, and conclusions are the sole responsibility of the
authors.

%% ----------------------------------------------------------------
%% CRediT authorship — prints the \credit{} entries above
%% ----------------------------------------------------------------
\printcredits

%% ================================================================
%% Bibliography
%% ================================================================
\bibliographystyle{cas-model2-names}
\bibliography{references}

%% Manual references — use until references.bib is ready
%%
%% \begin{thebibliography}{99}
%%
%% \bibitem[Baena~Neto and Becerra~Sablón(2026)]{BaenaNeto2026}
%% Baena~Neto, R., Becerra~Sablón, V.I., 2026. A speech analytics-based
%% methodological protocol for monitoring orthopedic rehabilitation in the
%% Brazilian Unified Health System.
%% Int. J. Environ. Res. Public Health 23, 626.
%% doi:10.3390/ijerph23050626
%%
%% \bibitem[Abbasian et~al.(2024)]{Abbasian2024}
%% Abbasian, M., et al., 2024. Foundation metrics for evaluating
%% effectiveness of healthcare conversations powered by generative AI.
%% npj Digit. Med. 7, 82.
%%
%% \bibitem[Alkhamees and Alasqah(2023)]{Alkhamees2023}
%% Alkhamees, M., Alasqah, I., 2023. Patient-physician communication in
%% intercultural settings: an integrative review.
%% Heliyon 9, e22667.
%%
%% \bibitem[Albert et~al.(2024)]{Albert2024}
%% Albert, P., et al., 2024. A scoping review of AI, speech and NLP methods
%% for assessment of clinician-patient communication. medRxiv.
%% doi:10.1101/2024.12.13.24318778 [preprint]
%%
%% \bibitem[Arksey and O'Malley(2005)]{Arksey2005}
%% Arksey, H., O'Malley, L., 2005. Scoping studies: towards a methodological
%% framework. Int. J. Soc. Res. Methodol. 8, 19--32.
%%
%% \bibitem[Bazoge et~al.(2023)]{Bazoge2023}
%% Bazoge, A., et al., 2023. Applying NLP to textual data from clinical
%% data warehouses: systematic review.
%% JMIR Med. Inform. 11, e42477.
%%
%% \bibitem[Fagherazzi et~al.(2021)]{Fagherazzi2021}
%% Fagherazzi, G., et al., 2021. Voice for health: the use of vocal
%% biomarkers from research to clinical practice.
%% Digit. Biomark. 5, 78--88.
%%
%% \bibitem[Imel et~al.(2022)]{Imel2022}
%% Imel, Z., et al., 2022. Identifying topics in patient and doctor
%% conversations using NLP methods. PCORI Final Research Report.
%% doi:10.25302/08.2021.me.160234167
%%
%% \bibitem[Kelley et~al.(2014)]{Kelley2014}
%% Kelley, J.M., et al., 2014. The influence of the patient-clinician
%% relationship on healthcare outcomes: a systematic review and
%% meta-analysis. PLoS ONE 9, e94207.
%%
%% \bibitem[Klusty et~al.(2025)]{Klusty2025}
%% Klusty, M.A., et al., 2025. Toward automated clinical transcriptions.
%% PMC12150720.
%%
%% \bibitem[Koleck et~al.(2019)]{Koleck2019}
%% Koleck, T.A., et al., 2019. NLP of symptoms documented in free-text
%% narratives of EHRs: a systematic review.
%% J. Am. Med. Inform. Assoc. 26, 364--379.
%%
%% \bibitem[Krishna et~al.(2021)]{Krishna2021}
%% Krishna, K., et al., 2021. Generating SOAP notes from doctor-patient
%% conversations. arXiv:2005.01795.
%%
%% \bibitem[Levac et~al.(2010)]{Levac2010}
%% Levac, D., et al., 2010. Scoping studies: advancing the methodology.
%% Implement. Sci. 5, 69.
%%
%% \bibitem[McCabe and Healey(2018)]{McCabe2018}
%% McCabe, R., Healey, P.G.T., 2018. Miscommunication in doctor-patient
%% communication. Top. Cogn. Sci. 10, 409--424.
%%
%% \bibitem[Muscat et~al.(2020)]{Muscat2020}
%% Muscat, D.M., et al., 2020. Health literacy and shared decision-making.
%% J. Gen. Intern. Med. 35, 2992--2999.
%%
%% \bibitem[Navarro et~al.(2023)]{Navarro2023}
%% Navarro, D.F., et al., 2023. Clinical NER and relation extraction using
%% NLP of medical free text: a systematic review.
%% Int. J. Med. Inform. 174, 105122.
%%
%% \bibitem[Peters et~al.(2020)]{Peters2020}
%% Peters, M.D.J., et al., 2020. Chapter 11: Scoping Reviews.
%% JBI Manual for Evidence Synthesis. JBI, Adelaide.
%%
%% \bibitem[Radford et~al.(2023)]{Radford2023}
%% Radford, A., et al., 2023. Robust speech recognition via large-scale
%% weak supervision. Proc. ICML 2023.
%%
%% \bibitem[Rucinski et~al.(2022)]{Rucinski2022}
%% Rucinski, K., et al., 2022. Patient and healthcare team adherence in
%% orthopaedics: a systematic review.
%% J. Patient Care 8, 203.
%%
%% \bibitem[Schloss and Konam(2020)]{Schloss2020}
%% Schloss, B., Konam, S., 2020. Towards an automated SOAP note.
%% Proc. Mach. Learn. Res. arXiv:2007.08749.
%%
%% \bibitem[Sezgin et~al.(2023)]{Sezgin2023}
%% Sezgin, E., et al., 2023. Extracting medical information from free-text
%% patient-generated health data using NLP. JMIR Form. Res. 7, e43014.
%%
%% \bibitem[Shao et~al.(2025)]{Shao2025}
%% Shao, J., et al., 2025. Development of a measurement of doctor-patient
%% communication quality scale.
%% Front. Public Health 13, 1606403.
%%
%% \bibitem[Tran et~al.(2023)]{Tran2023}
%% Tran, B.D., et al., 2023. ASR performance for digital scribes:
%% comparison between general-purpose and specialised models.
%% AMIA Annu. Symp. Proc. PMC10148344.
%%
%% \bibitem[Tricco et~al.(2018)]{Tricco2018}
%% Tricco, A.C., et al., 2018. PRISMA Extension for Scoping Reviews
%% (PRISMA-ScR). Ann. Intern. Med. 169, 467--473.
%%
%% \bibitem[vanBuchem et~al.(2021)]{vanBuchem2021}
%% van Buchem, M.M., et al., 2021. The digital scribe in clinical practice:
%% a scoping review and research agenda.
%% npj Digit. Med. 4, 57.
%%
%% \bibitem[Wang et~al.(2018)]{Wang2018}
%% Wang, Y., et al., 2018. Clinical information extraction applications:
%% a literature review.
%% J. Biomed. Inform. 77, 34--49.
%%
%% \bibitem[Watabe et~al.(2025)]{Watabe2025}
%% Watabe, S., et al., 2025. A patient-centered approach to developing
%% and validating an NLP model for extracting patient-reported symptoms.
%% Sci. Rep. 15. doi:10.1038/s41598-025-12845-3
%%
%% \end{thebibliography}

\end{document}